% !TeX root = ../main.tex
%------------------------------------------------------------------------------
%  Appendix D -- Closed-Form Solutions and Design Charts
%  Sources: Kirsch (1898); Hoek & Brown (1980) appendices; Carranza-Torres &
%  Fairhurst (2000); Muir Wood (1975); Duddeck & Erdmann (1985);
%  Einstein & Schwartz (1979); Panet (1995).
%------------------------------------------------------------------------------
\chapter{Closed-Form Solutions and Design Charts}
\label{app:closed-form}

A consolidated reference for the elastic and elasto-plastic solutions used
across Chapters~\ref{ch:ccm}, \ref{ch:ground-behaviour}, \ref{ch:openings}
and~\ref{ch:linings},
with their assumptions stated once so the main text need not repeat them.
Compression is positive throughout.

\section{Assumptions common to all solutions}
\label{sec:closed-form-assumptions}

Every solution in this appendix assumes a homogeneous, isotropic, linearly
elastic medium of infinite extent, plane strain conditions, and an opening
whose boundary is traction-free or carries a uniform internal pressure. Body
forces are neglected, so the stress field is uniform at infinity.

Each assumption fails in a characteristic way. \emph{Homogeneity} fails where
the opening crosses a boundary, and the solution then applies only to the
weaker unit. \emph{Isotropy} fails in bedded and foliated rock, where the
error is largest in the displacement field rather than the stresses.
\emph{Elasticity} fails wherever the strength is exceeded, which is precisely
where the answer matters --- hence the elasto-plastic solutions of
Section~\ref{sec:elasto-plastic-and-support-solutions}. \emph{Plane strain}
fails within about one diameter of the face
(Chapter~\ref{ch:rock-mechanics}). \emph{Infinite extent} fails at shallow
cover, conventionally below about two diameters, where the free surface
interacts with the opening and the solutions overestimate confinement.

\section{Circular and shaped openings}
\label{sec:circular-and-shaped-openings}

\subsection{Circular opening: the Kirsch solution}
\label{sec:kirsch-appendix}

The solution is that of \citet{kirsch1898}.
For a circular opening of radius $a$ in a biaxial field with vertical stress
$\pzero$ and horizontal stress $k\pzero$, the stresses at radius $r$ and angle
$\theta$ measured from the horizontal axis are

\begin{align}
\sigma_r &= \tfrac{\pzero}{2}\left[(1+k)\left(1-\tfrac{a^2}{r^2}\right)
   - (1-k)\left(1-\tfrac{4a^2}{r^2}+\tfrac{3a^4}{r^4}\right)\cos 2\theta\right],\\
\sigma_\theta &= \tfrac{\pzero}{2}\left[(1+k)\left(1+\tfrac{a^2}{r^2}\right)
   + (1-k)\left(1+\tfrac{3a^4}{r^4}\right)\cos 2\theta\right],\\
\tau_{r\theta} &= \tfrac{\pzero}{2}\left[(1-k)
   \left(1+\tfrac{2a^2}{r^2}-\tfrac{3a^4}{r^4}\right)\sin 2\theta\right].
\end{align}

At the boundary $r = a$ the radial and shear stresses vanish and the
tangential stress reduces to
\[
\sigma_\theta = \pzero\left[(1+k) + 2(1-k)\cos 2\theta\right],
\]
giving, at the sidewall ($\theta = 0$) and the crown
($\theta = 90^\circ$),
\[
\sigma_\theta^{\text{wall}} = \pzero\,(3-k), \qquad
\sigma_\theta^{\text{crown}} = \pzero\,(3k-1).
\]

Two results follow directly. In a hydrostatic field ($k=1$) the boundary
stress is $2\pzero$ everywhere. And the crown goes into tension whenever
$k < 1/3$, which is the origin of the destressed crown discussed in
Chapter~\ref{ch:ground-behaviour}.

The radial displacement at the boundary of an unsupported opening in a
hydrostatic field is
\[
u_r = \frac{\pzero\,a\,(1+\nu)}{E}.
\]

\subsection{Boundary stress concentration by opening shape}
\label{sec:shape-factors}

Boundary stresses for non-circular openings are conveniently expressed as
concentration factors applied to the field stress. Table~\ref{tab:shape-factors}
gives crown and sidewall factors for circular and elliptical profiles in a
hydrostatic field, from the expressions of Section~\ref{sec:elliptical}, and
comparative indications for angular profiles; for other $k$ the values must
be obtained from the elliptical solution or numerically.

\begin{table}[htbp]
\caption{Boundary stress concentration by profile, hydrostatic field.
Corner values, not crown or sidewall values, govern angular profiles.}
\label{tab:shape-factors}
\centering
\small
\begin{tabularx}{\textwidth}{l c c X}
\toprule
\textbf{Profile} & \textbf{Crown} & \textbf{Sidewall} & \textbf{Remarks} \\
\midrule
Circular & 2.0 & 2.0 & Uniform in a hydrostatic field \\
Elliptical, wide ($W/H = 2$) & 1.0 & 4.0 & Crown $2H/W$, sidewall $2W/H$; suited to $k > 1$ \\
Elliptical, tall ($W/H = 1/2$) & 4.0 & 1.0 & The same expressions; suited to $k < 1$ \\
Horseshoe & moderate & higher & Sharp invert corners require radiusing \\
Square, radiused corners & low & low & Corner concentration governs \\
\bottomrule
\end{tabularx}
\end{table}

For angular profiles the governing value is at the corners, where the
concentration rises steeply as the radius decreases. This is why square
profiles are radiused and why a flat roof is avoided in high horizontal
stress.

\subsection{Elliptical openings and the optimum shape}
\label{sec:elliptical}

For an ellipse with semi-axis $W$ parallel to the horizontal stress and $H$
parallel to the vertical, in a field with vertical stress $\pzero$ and
horizontal $k\pzero$, the boundary stresses at the ends of the axes are

\[
\sigma_{\text{wall}} = \pzero\left(2\frac{W}{H} + 1 - k\right),
\qquad
\sigma_{\text{crown}} = \pzero\left(k\left(2\frac{H}{W} + 1\right) - 1\right).
\]

Setting the two equal gives the \emph{harmonic hole} condition, at which the
boundary stress is uniform around the opening and the peak concentration is
minimised:
\[
\frac{W}{H} = k .
\]

The optimum opening is therefore elongated in the direction of the major
principal stress, in the ratio of the stresses. A cavern in a field with
$k = 2$ is best made twice as wide as it is high, when crown and wall both
carry $3\pzero$; the same cavern rotated through ninety degrees carries a
crown stress of $9\pzero$, three times the optimum value.

\section{Elasto-plastic and support solutions}
\label{sec:elasto-plastic-and-support-solutions}

\subsection{Elasto-plastic solutions}
\label{sec:elasto-plastic-appendix}

For a circular opening of radius $a$ in a hydrostatic field $\pzero$ with
internal support pressure $\pin$, in Mohr--Coulomb ground with cohesion $c$
and friction angle $\phi$, define
\[
k_\phi = \frac{1+\sin\phi}{1-\sin\phi},
\qquad
\sigma_{cm,MC} = \frac{2c\cos\phi}{1-\sin\phi}.
\]

Yielding begins when the support pressure falls below the critical value
\[
p_{cr} = \frac{2\pzero - \sigma_{cm,MC}}{1 + k_\phi}.
\]

For $\pin < p_{cr}$ the plastic zone extends to radius
\[
R_p = a\left[\frac{2\left(\pzero(k_\phi-1) + \sigma_{cm,MC}\right)}
{(1+k_\phi)\left((k_\phi-1)\pin + \sigma_{cm,MC}\right)}\right]^{\frac{1}{k_\phi-1}} .
\]

\noindent\textbf{Which strength.}\enspace These two expressions take the
unconfined compressive strength of a Mohr--Coulomb material, which is why it is
written $\sigma_{cm,MC}$ here. The bare symbol $\scm$ is the Hoek--Brown rock
mass strength, in the List of Symbols and in
Appendix~\ref{app:calculations}, and the two are not interchangeable. In the
worked example of Chapter~\ref{ch:ccm} --- $\sci = 72$\,MPa, $\GSI = 30$,
$\mi = 7$, $\pzero = 16.9$\,MPa, with the Mohr--Coulomb fit taken to
$\sigma_3 \approx \pzero$ --- they are about 5.3 and 1.2\,MPa, and substituting
the Hoek--Brown value for the Mohr--Coulomb one inflates the plastic radius by
about two and a half times and the wall convergence by about six. Both factors
grow with $\pzero$ and neither is a property of strength alone. Where a
sequence in Appendix~\ref{app:calculations} delivers a rock mass strength, $c$
and $\phi$ are fitted first and $\sigma_{cm,MC}$ formed from them.

\noindent\textbf{Use.}\enspace The critical pressure is the more useful of the
two results, because it answers immediately whether the ground will yield at
all. Where $p_{cr} \leq 0$ the opening is self-supporting in the
elastic sense, and the design problem is structural rather than one of ground
pressure.

Equivalent expressions for Hoek--Brown ground follow the formulation of
\citet{carranza2000ccm}, in which the same quantities are obtained through a
scaled stress transformation. They are algebraically longer and identical in
structure and use.

\subsection{Ground reaction and support characteristic expressions}
\label{sec:grc-scc-appendix}

The ground reaction curve relates support pressure to wall convergence. In the
elastic range,
\[
u_r = \frac{(\pzero - \pin)\,a\,(1+\nu)}{E},
\]
and beyond yield it is obtained from the plastic solution above, with the
convergence increasing rapidly as $\pin$ falls. For a material that yields
without volume change,
\[
u_r \approx u_r(p_{cr})\left(\frac{R_p}{a}\right)^{2} ,
\]
where $u_r(p_{cr})$ is the elastic convergence at the onset of yield; a
dilatant flow rule increases it. This is the form used in
Chapter~\ref{ch:ccm}.

Support stiffness $k_s$ and maximum pressure $p_{sm}$ for the common
systems are as follows, with $a$ the tunnel radius throughout.

\emph{Shotcrete or concrete ring} of thickness $t_c$, modulus $E_c$ and
Poisson's ratio $\nu_c$, with unconfined strength $\sigma_{cc}$:
\[
k_s = \frac{E_c\left[a^2-(a-t_c)^2\right]}
{(1+\nu_c)\left[(1-2\nu_c)a^2+(a-t_c)^2\right]a},
\qquad
p_{sm} = \frac{\sigma_{cc}}{2}\left[1-\frac{(a-t_c)^2}{a^2}\right].
\]

\emph{Segmental rings} are represented under the uniform convergence assumed
here by the same expressions, with $t_c$ the segment thickness and $E_c$,
$\nu_c$ and $\sigma_{cc}$ those of the segment concrete. Under axisymmetric
load the radial joints carry hoop thrust in compression without rotating, so
the stiffness is close to that of a continuous ring; the joints matter under
non-uniform load, which is the province of the lining solutions of
Section~\ref{sec:charts-and-lining-solutions} and of
Chapter~\ref{ch:linings}. Two quantities differ from those of a cast ring.
The capacity is governed at the radial joints, where the hoop thrust $\pin a$
passes through a contact band of height $h_j$, less than $t_c$
(Section~\ref{sec:joint-behaviour}), so that
\[
p_{sm} = \frac{\sigma_{cc}\,h_j}{a} ,
\]
the thin-ring form with $h_j$ in place of $t_c$. And the convergence $u_{s0}$
at which the ring begins to carry load is fixed by the tail void and the
effectiveness of the annulus grout, not read from the longitudinal
deformation profile.

\emph{Steel sets} of cross-sectional area $A_s$, modulus $E_s$ and yield
stress $\sigma_{ys}$, at longitudinal spacing $s$:
\[
k_s = \frac{E_s A_s}{s\,a^2},
\qquad
p_{sm} = \frac{A_s \sigma_{ys}}{s\,a}.
\]
Both must be reduced to allow for block seating. The blocking between set and
rock deforms in series with the set, so that the effective stiffness is
$\left(1/k_{s,\text{set}} + 1/k_{s,\text{block}}\right)^{-1}$ and the softer
element governs; a poorly seated set can therefore fall well below the
computed values. No general factor applies, and the reduction must be
estimated from the blocking detail --- the number, spacing and stiffness of
the blocks.

\emph{Mechanically or chemically anchored bolts and cables} (ungrouted) of free length $L_b$, diameter $d_b$
and modulus $E_b$, at spacing $s_c$ circumferentially and $s_l$
longitudinally, with a load--deformation constant $Q$ for the anchor and
plate:
\[
k_s = \frac{1}{s_c s_l}\left[\frac{4L_b}{\pi d_b^2 E_b} + Q\right]^{-1},
\qquad
p_{sm} = \frac{T_{bf}}{s_c s_l},
\]
where $T_{bf}$ is the bolt capacity. The expression treats the bar as an elastic tie between anchor and plate; a fully grouted bar is not a tie but reinforcement of the rock mass, and for it the expression is only an approximation. The constant $Q$, in metres per
meganewton, is read from the load--extension curve of a pull-out test on the
anchor and plate assembly. It depends on the anchor type, the hole and the
rock, no representative range is tabulated in Appendix~\ref{app:constants},
and it is measured on site before the expression is used; representative
values of $E_b$ and of the steel-set properties are given there.

For systems installed together, stiffnesses add:
$k_{s,\text{total}} = \sum k_{s,i}$, and the combined capacity is reached when the
first element attains its own $p_{sm}$ --- not the sum of the individual
capacities.

\subsection{Longitudinal deformation profile expressions}
\label{sec:ldp-appendix}

The longitudinal deformation profile gives the convergence at distance $X$
behind the face as a fraction of the final convergence $u_{\max}$. The widely
used elastic fit of \citet{panet1995ccm} is
\[
\frac{u(X)}{u_{\max}} = 0.25 + 0.75\left[1-\left(\frac{0.75}{0.75+X/a}\right)^{2}\right],
\]
which gives 25\,\% of the total convergence at the face, about 90\,\%
within one and a third radii behind it and about 97\,\% at three radii.

For elasto-plastic ground the profile depends on the extent of the plastic
zone, and the profile of \citet{vlachopoulos2009ldp}, written in terms of
$R_p/a$, is used (Section~\ref{sec:scc-and-ldp}). In yielding ground the
profile is drawn out over a greater distance, and the elastic fit overstates
the fraction of the convergence already released at a given distance behind
the face. It therefore understates the load the support attracts, which is
unconservative for support design.

\noindent\textbf{Use.}\enspace The profile converts the distance behind the
face at which support is installed into the convergence already released,
which is the offset applied to the support characteristic curve in the
verification sequence of Appendix~\ref{app:calculations}.

The plastic solution, the ground reaction curve and the combined support
characteristic of this section are evaluated in \companionname\
(Appendix~\ref{app:software}), which also returns the compressibility and
flexibility ratios of Section~\ref{sec:charts-and-lining-solutions}.

\section{Charts and lining solutions}
\label{sec:charts-and-lining-solutions}

\paragraph{Zone of influence.}
The zone of influence is conventionally bounded where the stress disturbance
falls below 5\,\% of the field stress. For a circular opening in a hydrostatic
field, the Kirsch solution gives the disturbance in $\sigma_\theta$ as
$(a/r)^2$, so the 5\,\% contour lies at
\[
r \approx 4.5\,a ,
\]
that is, about two and a quarter diameters from the centre, or one and
three-quarter diameters clear of the boundary. For non-hydrostatic fields the
shape depends on the stress to which the 5\,\% is referred. Referred to
$\pzero$ on both axes, the zone is elongated modestly in the direction of the
minor principal stress: at $k = 0.5$ the contour lies at about $4.0\,a$ along
the horizontal axis and $3.7\,a$ along the vertical. Referred instead to the
far-field tangential stress on each axis, $\pzero$ on the horizontal and
$k\pzero$ on the vertical, it lies at about $4.0\,a$ and $5.4\,a$, elongated
towards the major principal stress. The convention should therefore be stated
with the result.

Two openings whose zones of influence do not overlap may be analysed
independently. Closer, they interact, and the treatment is
Chapter~\ref{ch:openings}.

\paragraph{Lining solutions.}
For a circular lining of cross-sectional area $A_l$ and second moment $I_l$,
both per unit length of tunnel, with modulus $E_l$ and Poisson's ratio
$\nu_l$, in ground of modulus $E$ and Poisson's ratio $\nu$, define the
dimensionless compressibility and flexibility ratios
\[
C = \frac{E\,a\,(1-\nu_l^2)}{E_l A_l (1-\nu^2)},
\qquad
F = \frac{E\,a^3(1-\nu_l^2)}{E_l I_l (1-\nu^2)} .
\]

The thrust and moment then follow as coefficients of the field stress. For the
full-slip condition and a lining installed before the ground deforms, the
solution of \citet{einstein1979simplified} is, with $\theta$ measured from the
springline as in the Kirsch solution,
\begin{align*}
T(\theta) &= \tfrac{1}{2}\pzero a\left[(1+k)(1-a_0^{*})
  + (1-k)(1-2a_2^{*})\cos 2\theta\right],\\
M(\theta) &= \tfrac{1}{2}\pzero a^2 (1-k)(1-2a_2^{*})\cos 2\theta,
\end{align*}
\[
a_0^{*} = \frac{C F (1-\nu)}{C + F + C F (1-\nu)},
\qquad
a_2^{*} = \frac{(F+6)(1-\nu)}{2F(1-\nu) + 6(5-6\nu)},
\]
where the thrust $T$ is compressive, and the coefficients $a_0^{*}$ and
$a_2^{*}$, which are not the radius $a$, are dimensionless. The moments at crown
and springline are equal and opposite, the crown lining being in tension on its
inner face when $k < 1$. The formulations of \citet{muirwood1975} and
\citet{duddeck1985structural} differ in detail and agree in behaviour.

\noindent\textbf{What the ratios mean.}\enspace The ratio $C$ governs how much
of the ground load appears as thrust, and $F$ how much appears as moment. At
high $F$ the lining is flexible relative to the ground: it conforms, and the
moment tends to zero, the thrust being set by $C$. At low $F$ the lining
resists, distorts less, and attracts large moments. Since a lining designed
for thrust alone is far more efficient than one designed for moment, high $F$
is generally the better design --- the result quoted in
Chapter~\ref{ch:linings}.

Full-slip and no-slip solutions bracket the real interface condition; the
no-slip case gives higher thrust and lower moment, and both should be checked.

\section{Limits of use}
\label{sec:closed-form-limits}

These solutions remain the right tool for preliminary sizing, for
understanding which variables govern, and --- most valuably --- as a check on
numerical output. A numerical model of a circular opening in uniform ground
that does not reproduce the Kirsch stresses to within a few per cent contains
an error that will not announce itself in the complex geometry the model was
built for.

They serve only as a check where the geometry is non-circular and the ground
non-uniform, where construction is staged, or where openings interact. They
should not be used at all at shallow cover below about two diameters, in
strongly anisotropic rock, where the excavation is not remote from the face
--- the longitudinal deformation profile of Section~\ref{sec:ldp-appendix}
excepted, being an empirical fit to three-dimensional results for precisely
that region --- or where the ground is time-dependent. In each of those cases the assumption
violated is not a refinement but the basis of the derivation.

\chapterreferences
